% Problem record op_161b5f34e4751e33. This TeX file is derived from
% database/problems_json/op_161b5f34e4751e33.json by scripts/sync-tex.mjs; edit the JSON record
% and rerun the script rather than editing this file.
\section{Polynomial fractional coloring for Pauli shadow tomography}
\label{sec:op_161b5f34e4751e33}

\paragraph{Problem.}
Shadow tomography asks how to estimate many expectation values using copies of an unknown state. On $n$ qubits, a Pauli observable is a tensor product $P\in\{I,X,Y,Z\}^{\otimes n}$. Its expectation on a state $\rho$ is $\operatorname{Tr}(\rho P)$. For an accuracy scale $0<\varepsilon<1$, let $B_\varepsilon(\rho)$ be the set of Paulis with $\lvert\operatorname{Tr}(\rho P)\rvert\geq\varepsilon$. The question is how economically these observables can be grouped for measurement.

Build a graph with vertex set $B_\varepsilon(\rho)$ and join $P$ to $Q$ when they anticommute, meaning $PQ=-QP$. An independent set $S$ is a group of pairwise commuting Paulis. The observables in one such group have a common eigenbasis, so a single joint measurement yields an outcome for each of them.

A fractional coloring assigns a weight $w_S\geq0$ to each independent set. Every vertex $P$ must be covered to total weight at least one, meaning $\sum_{S\ni P}w_S\geq1$. Write $\chi_f$ for the minimum total weight $\sum_S w_S$ of such a cover. A cover of weight $W$ can be used as a measurement schedule. Choosing $S$ with probability $w_S/W$ includes each retained Pauli with probability at least $1/W$.

Do universal constants $C,\kappa>0$ exist such that $\chi_f\leq C\varepsilon^{-\kappa}$ for every number of qubits $n$, every state $\rho$, and every $0<\varepsilon<1$? The weights may depend on the state. The question asks for the existence of the cover.

\subsection*{Status}

Unsolved

\subsection*{Source}

Stempin, Llorens, and Huber pose the polynomial bound in Sec.~7 after disproving the quadratic conjecture \sourcecite{ref:pauli-coloring}{SLH26}.

\subsection*{Progress}

\begin{itemize}
  \item The anti-heptagon construction in Sec.~5 and its generalization in Sec.~6 give families with unbounded $\varepsilon^2\chi_f$. The polynomial-bound question allows a larger exponent \sourcecite{ref:pauli-coloring}{SLH26}.

  \item King et al.\ relate fractional colorings to protocols for estimating Pauli expectations using pairs of copies. Efficient construction or sampling of the coloring is an additional computational requirement \sourcecite{ref:triply-efficient}{KGKB25}.
\end{itemize}

\subsection*{References}

\begin{enumerate}[label={},leftmargin=4.8em,labelsep=0.5em]
  \footnotesize
  \item[\textup{[SLH26]}]\label{ref:pauli-coloring}
  J. Stempin, S. Llorens, and F. Huber, \emph{Counterexamples to the fractional coloring conjecture for triply efficient shadow tomography} (2026). \href{https://arxiv.org/abs/2608.20113v1}{arXiv:2608.20113v1}.

  \item[\textup{[KGKB25]}]\label{ref:triply-efficient}
  R. King, D. Gosset, R. Kothari, and R. Babbush, \emph{Triply Efficient Shadow Tomography}, PRX Quantum \textbf{6}, 010336 (2025). \href{https://doi.org/10.1103/PRXQuantum.6.010336}{DOI:10.1103/PRXQuantum.6.010336}.
\end{enumerate}

\subsection*{Comment}

The polynomial bound would give a dimension-independent cover cost for the relevant commuting measurement groups. Computing the cover is a separate question. The catalog and literature were checked through November 26, 2026.

\subsection*{Field}

Quantum Statistical Inference

\subsection*{Topic}

Quantum tomography

\subsection*{ID}

\texttt{op\_161b5f34e4751e33}
